This thesis focuses on the modeling and control of a plate heat exchanger using predictive control methods. The thesis aimed to identify system dynamics from experimental measurements, to design Model Predictive Control (MPC) controllers using two modeling approaches, and to compare their performance in simulations and experiments.
System identification was carried out using step changes in the power of the hot fluid pump. The pump power was used as the system input, while the temperature on the outlet from the heat exchanger was the system output. Based on the measured data, a first-order linear Strejc model was identified. In addition, a data-driven Koopman model was developed, which allows improved approximation of nonlinear system behavior using linear dynamics in a higher-dimensional space.
For both models, MPC controllers were designed with a quadratic optimization problem and with constraints on the system input and output. A Kalman filter was included in the control structure to estimate the system states. The models were compared in terms of control accuracy, deviation from the steady state, and control effort. The simulation results were complemented by a comparison with a nonlinear mathematical model of the heat exchanger.
In the experimental part of the thesis, MPC control was applied to the heat exchanger, where the system was controlled from different initial temperatures to a steady operating point. The obtained results confirmed the Koopman-based MPC achieved better results with comparable control effort.
The thesis demonstrates the applicability of the Koopman operator for the control of nonlinear thermal intensive systems and highlights its advantages compared to classical linear models such as the Strejc model.
